\documentclass[11pt,a4paper]{article}

\usepackage[T1]{fontenc}
\usepackage[utf8]{inputenc}
\usepackage[french]{babel}
\usepackage{lmodern}
\usepackage{geometry}
\usepackage{xcolor}
\usepackage{hyperref}
\usepackage{listings}
\usepackage{enumitem}
\usepackage{booktabs}
\usepackage{csquotes}
\usepackage{seqsplit}



\usepackage{array}
\usepackage{ragged2e}

\geometry{margin=2.5cm}
\hypersetup{
  colorlinks=true,
  linkcolor=blue!50!black,
  urlcolor=blue!50!black,
  pdftitle={Guide d'utilisation - LINGO vers Pyomo},
  pdfauthor={Projet lingo_to_pyomo}
}

\definecolor{codebg}{RGB}{248,248,248}
\definecolor{codeframe}{RGB}{220,220,220}
\definecolor{codecomment}{RGB}{0,120,0}
\definecolor{codekeyword}{RGB}{0,70,170}

\lstdefinestyle{mystyle}{
  backgroundcolor=\color{codebg},
  frame=single,
  rulecolor=\color{codeframe},
  basicstyle=\ttfamily\small,
  keywordstyle=\color{codekeyword}\bfseries,
  commentstyle=\color{codecomment},
  showstringspaces=false,
  breaklines=true,
  columns=fullflexible
}

\lstset{style=mystyle}

\title{Guide d'utilisation\\\large Conversion LINGO $\rightarrow$ Pyomo}

\begin{document}
\maketitle

\section{Objectif du projet}
Ce projet convertit des modèles d'optimisation écrits en LINGO (fichiers \texttt{.lng}) vers du code Pyomo, puis génère un notebook Jupyter prêt à être exécuté.

\medskip
\noindent
Le point important est le suivant :
\begin{itemize}[leftmargin=1.2cm]
  \item vous pouvez l'utiliser facilement via l'application web Flask ;
  \item vous pouvez aussi l'utiliser directement comme une bibliothèque Python classique, fonction par fonction.
\end{itemize}

\section{Installation rapide}
\begin{lstlisting}[language=bash]
git clone https://github.com/Joaquim2805/lingo_to_pyomo.git
cd lingo_to_pyomo
python -m venv .venv
source .venv/bin/activate
pip install -r requirements.txt
\end{lstlisting}

\noindent
Ensuite, placez vos fichiers LINGO dans le dossier \texttt{data/} (ou utilisez ceux déjà présents dans le projet).

\section{Deux façons d'utiliser l'outil}
\subsection{A. Via l'application web Flask (parcours conseillé)}
Cette voie est la plus simple pour un usage quotidien.

\begin{lstlisting}[language=bash]
cd dev
python app.py
\end{lstlisting}

Puis ouvrez \texttt{http://localhost:5000} dans le navigateur.

\medskip
\noindent
Depuis l'interface, vous avez 4 actions principales :
\begin{enumerate}[leftmargin=1.2cm]
  \item \textbf{Générer Notebook} : conversion directe LINGO $\rightarrow$ Pyomo.
  \item \textbf{Conversion OLE} : permet de convertir un fichier LINGO utilisant un fichier Excel comme source de données externes via \texttt{@OLE(...)} en fichier LINGO explicite.
  \item \textbf{Nettoyage} : normalise un fichier LINGO (commentaires, accents, casse, sections \texttt{DATA/ENDDATA} vides).
  \item \textbf{Pipeline complet} : enchaîne automatiquement les étapes nécessaires, avec nettoyage optionnel, génération du notebook, et export JSON optionnel des données.
\end{enumerate}



\subsection{B. Comme bibliothèque Python (mode script/notebook)}
Ce mode est pratique pour l'automatisation et les tests.

Le pipeline standard est :
\begin{enumerate}[leftmargin=1.2cm]
  \item parser le fichier LINGO ;
  \item transformer l'arbre en structure Python ;
  \item générer le code Pyomo ;
  \item générer un notebook Jupyter ;
  \item (optionnel) externaliser les données en JSON.
\end{enumerate}

\begin{lstlisting}[language=Python]
from src.lingo_parser.parser import parse_lingo_model
from src.lingo_parser.transformer import LingoModelTransformer2
from src.pyomo_generator.json_parser import generate_pyomo_code
from src.notebook_generator.notebook_construct import generate_pyomo_notebook

tree = parse_lingo_model("data/Pastissimo.lng")
model_dict = LingoModelTransformer2().transform(tree)
pyomo_code = generate_pyomo_code(model_dict, external_data=False)
generate_pyomo_notebook(pyomo_code, solver="highs", filename="notebooks/Pastissimo.ipynb")
\end{lstlisting}

\noindent
Le notebook de développement du projet montre exactement ce style d'utilisation (appels directs aux fonctions dans une cellule Python).

\section{Exemples concrets}
\subsection{Exemple 1 : Conversion simple d'un fichier \texttt{.lng}}
\begin{lstlisting}[language=Python]
from src.lingo_parser.parser import parse_lingo_model
from src.lingo_parser.transformer import LingoModelTransformer2
from src.pyomo_generator.json_parser import generate_pyomo_code
from src.notebook_generator.notebook_construct import generate_pyomo_notebook

tree = parse_lingo_model("data/Buckly.lng")
model_dict = LingoModelTransformer2().transform(tree)
pyomo_code = generate_pyomo_code(model_dict, external_data=False)

generate_pyomo_notebook(
    pyomo_code,
    solver="highs",
    filename="notebooks/Buckly.ipynb"
)
\end{lstlisting}

\subsection{Exemple 2 : Données externes JSON}
Externaliser les données en JSON signifie séparer le \textit{modèle} (équations, contraintes, variables) des \textit{valeurs de données} (paramètres, tableaux, coefficients).
Le code Pyomo reste alors stable, et les données sont lues depuis un fichier \texttt{.json} au moment de l'exécution. Cette approche est utile pour tester plusieurs jeux de données sans régénérer tout le code.

\begin{lstlisting}[language=Python]
from src.lingo_parser.parser import parse_lingo_model
from src.lingo_parser.transformer import LingoModelTransformer2
from src.pyomo_generator.json_parser import generate_pyomo_code, save_pyomo_data_to_json
from src.notebook_generator.notebook_construct import generate_pyomo_notebook

tree = parse_lingo_model("data/Cardoza.lng")
model_dict = LingoModelTransformer2().transform(tree)

# Sauvegarde des données du modèle
save_pyomo_data_to_json(model_dict, output_path="data/Cardoza_data.json")

# Code Pyomo qui lit les données externes
pyomo_code = generate_pyomo_code(
  model_dict,
  external_data=True,
  data_filename="../data/Cardoza_data.json"
)

generate_pyomo_notebook(
  pyomo_code,
  solver="highs",
  filename="notebooks/Cardoza.ipynb",
  external_data=True,
  json_data_filename="../data/Cardoza_data.json"
)
\end{lstlisting}

\noindent
\subsection{Exemple 3 : Script utilitaire réutilisable}
\begin{lstlisting}[language=Python]
from src.lingo_parser.parser import parse_lingo_model
from src.lingo_parser.transformer import LingoModelTransformer2
from src.pyomo_generator.json_parser import generate_pyomo_code
from src.notebook_generator.notebook_construct import generate_pyomo_notebook

def convertir_modele(input_lng, output_ipynb, solver="highs"):
  tree = parse_lingo_model(input_lng)
  model_dict = LingoModelTransformer2().transform(tree)
  pyomo_code = generate_pyomo_code(model_dict, external_data=False)
  generate_pyomo_notebook(pyomo_code, solver=solver, filename=output_ipynb)


convertir_modele("data/Pastissimo.lng", "notebooks/Pastissimo.ipynb")
\end{lstlisting}

\section{Fonctions principales à connaître}
\begin{description}
    \item[\texttt{parse\_lingo\_model(...)}] Lit un fichier LINGO et renvoie un arbre syntaxique (Lark).
    
    \item[\texttt{LingoModelTransformer2().transform(tree)}] Convertit l'arbre syntaxique en dictionnaire Python exploitable.
    
    \item[\texttt{generate\_pyomo\_code(...)}] Produit le code Pyomo sous forme de chaîne de caractères.
    
    \item[\texttt{generate\_pyomo\_notebook(...)}] Construit le notebook \texttt{.ipynb}.
    
    \item[\texttt{save\_pyomo\_data\_to\_json(...)}] Exporte les données du modèle en JSON.
    
    \item[\texttt{clean\_lingo\_content(...)} / \texttt{clean\_lingo\_file(...)}] Nettoie le contenu LINGO.
\end{description}

\section{Que contient le notebook généré ?}
Le notebook généré est directement exploitable : vous l'ouvrez, vous exécutez les cellules, et vous obtenez la solution.

\medskip
\noindent
Concrètement, vous y trouverez :
\begin{enumerate}[leftmargin=1.2cm]
  \item les imports nécessaires (Pyomo, pandas, etc.) ;
  \item le modèle complet reconstruit en Python (sets, paramètres, variables, contraintes, objectif) ;
  \item une cellule qui lance le solveur choisi ;
  \item l'état de la résolution (faisable, optimal, etc.) ;
  \item un affichage des variables de décision non nulles, pour lire rapidement la solution utile.
\end{enumerate}




\section{Résolution et solveurs}
Dans les notebooks générés, la résolution passe par :
\begin{lstlisting}[language=Python]
solver = SolverFactory("highs")
result = solver.solve(model, tee=True)
\end{lstlisting}

\noindent
Les formulaires web proposent notamment : \texttt{gurobi}, \texttt{cplex}, \texttt{glpk}, \texttt{highs} (et \texttt{ipopt} dans le pipeline complet).

\section{Conclusion}
Ce projet peut être utilisé de deux manières complémentaires :
\begin{itemize}[leftmargin=1.2cm]
  \item \textbf{interface web Flask} pour un usage rapide et pédagogique ;
  \item \textbf{API Python} pour intégrer la conversion dans des scripts ou notebooks personnels.
\end{itemize}

\noindent
Dans les deux cas, le même cœur de conversion est utilisé : parsing LINGO, transformation intermédiaire, génération Pyomo, puis génération du notebook.

\end{document}
